\documentclass[reprint,superscriptaddress,amsmath,amssymb,aps,prd,showkeys,nofootinbib]{revtex4-2}

\usepackage{graphicx}
\usepackage{dcolumn}
\usepackage{bm}
\usepackage{hyperref}
\usepackage{microtype}
\usepackage{booktabs}

\begin{document}

\title{General-Relativistic Cyclic Cosmology from Holographic Saturation}

\author{Paul Jarvis}
\email{mrpaulwjarvis@gmail.com}
\affiliation{Independent Researcher, United Kingdom}
\date{\today}

\begin{abstract}
We present a cyclic cosmological framework in which the big-bang singularity is replaced by
an information-saturation event that triggers a transition from gravitational collapse to
expansion. The mechanism is formulated entirely within general relativity at the macroscopic
level: the transition is implemented as an Israel--Darmois junction across a spacelike
hypersurface at a critical density, interpreted as a causal boundary rather than a topology
change. The key dynamical ingredient is an effective confinement potential that reverses sign
at holographic saturation, converting net attraction into repulsion. Entropy is inherited
across cycles but renormalised relative to the enlarged phase space activated at the throat,
yielding an effectively low-entropy initial state without entropy destruction. Anisotropies
are suppressed by chaotic mixing in the collapsing phase and holographic coarse-graining at
saturation, producing a homogeneous FLRW initial geometry. We demonstrate numerically that
the confinement--rupture map defines a renormalisation-group-like flow on the space of
effective couplings with a stable fixed point, driving the renormalised inhomogeneity ratio
exponentially to zero across successive cycles. New results presented here establish the
precise spectral foundations of the underlying substrate: the MI-Laplacian spectral invariant
$\rho_0$ is identified as the integrated spectral weight at the hyper-stiff transition
($\alpha=6$, $w_{\rm eff}=1$), the tensor spectral index $n_T \simeq 1.875$ is derived from
full Mukhanov--Sasaki integration through a finite-duration transient phase (resolving the
discrepancy with the eternal-$w$ analytic formula), and the distinction between the Weyl
spectral dimension $d_{\rm eff}^{\rm Weyl}$ and the fluid dimension $d_{\rm fluid}=3(1+w)$
is established rigorously. The GW amplitude is shown to depend on the microscopic substrate
energy scale $\varepsilon_0$, while the spectral tilt and UV behaviour are parameter-independent
predictions of the framework.
\end{abstract}

\keywords{Cyclic Cosmology, Holographic Principle, Quantum Entanglement, Bouncing Universes,
Primordial Gravitational Waves}

\maketitle

% ─────────────────────────────────────────────────────────────────────────────
\section{Introduction}
% ─────────────────────────────────────────────────────────────────────────────

Classical general relativity (GR) predicts geodesic incompleteness under broad conditions,
as established by the Penrose--Hawking singularity theorems \cite{Penrose1965,
HawkingPenrose1970}. This motivates the search for frameworks in which the big-bang
singularity is replaced by a well-defined transition. Existing approaches include loop
quantum cosmology (LQC) \cite{AshtekarBojowald2006}, conformal cyclic cosmology (CCC)
\cite{PenroseBook}, ekpyrotic models \cite{SteinhardtTurok2002}, and holographic or
fuzzball-inspired constructions \cite{tHooft1993, Susskind1995}.

We develop a cyclic cosmology in which the singularity is replaced by an
information-saturation event. The transition from collapse to expansion is implemented as a
GR-consistent junction across a spacelike hypersurface at a critical density. The underlying
intuition is holographic: as a collapsing region approaches its saturation radius, its
information content approaches the Bekenstein--Hawking bound, and the effective confinement
potential undergoes a sign reversal.

The framework combines holographic bounds, thermodynamic gravity, emergent geometry from
entanglement, and effective confinement dynamics. It remains fully GR-compatible at the
macroscopic level while connecting to a pre-geometric substrate supported by robust
numerical evidence. The present work extends earlier results \cite{JarvisEEG2026,
JarvisGW2026} by establishing the precise spectral foundations of the substrate energy
density, resolving the tensor spectral index, and rigorously separating the Weyl spectral
dimension from the fluid equation-of-state dimension.

% ─────────────────────────────────────────────────────────────────────────────
\section{Pre-Geometric Substrate and Emergent Geometry}
% ─────────────────────────────────────────────────────────────────────────────

\subsection{Substrate Assumptions}
We assume that smooth geometry emerges natively from the quantum entanglement structure of
an underlying pre-geometric many-body substrate \cite{VanRaamsdonk2010}. The mutual
information between two arbitrary sub-regions $i$ and $j$,
\begin{equation}
I(i,j) = S(\rho_i) + S(\rho_j) - S(\rho_{ij}),
\end{equation}
serves as a relational metric proxy that maps spatial proximity. This mutual information
matrix dictates a symmetric connectivity matrix $A_{ij}$, from which we construct the
combinatorial graph Laplacian
\begin{equation}
L^{\rm MI}_{ij} = \delta_{ij}\sum_k I_{ik} - I_{ij}.
\end{equation}
No specific microscopic Hamiltonian or prior spatial embedding is assumed. All macroscopic
phenomenological parameters are treated as coarse-grained descriptions whose exact
microphysical derivations are deferred to future work.

\subsection{Heisenberg-Chain and Multi-Dimensional Spectral Maps}
To establish the validity of this mechanism, we track the low-lying spectrum of the graph
Laplacian derived from explicit quantum states. Initial benchmarks on the ground state of an
antiferromagnetic Heisenberg spin chain \cite{JarvisCode},
\begin{equation}
H = J\sum_{i=1}^{N-1}\vec{S}_i\cdot\vec{S}_{i+1},\quad J>0,
\end{equation}
for chain lengths $N=6$--$16$ confirm the emergence of a smooth one-dimensional geometry,
with a quadratic dispersion relation and a finite-size scaling law $\lambda_1\propto N^{-2}$.

Numerical diagonalisation of a relational entanglement kernel on a discrete 2D toroidal
lattice demonstrates spontaneous coordinate chart generation. As shown in
Fig.~\ref{fig:eigenmodes}, the lowest non-trivial modes reorganise into smooth standing
waves. The first excited tier contains a degenerate pair at $\lambda=0.675$, providing
the orthogonal spatial coordinate gradients required to assemble a continuous, isotropic
manifold.

\begin{figure}[t]
\includegraphics[width=\linewidth]{eigenmodes.png}
\caption{Emergent 2D eigenmodes from the substrate entanglement kernel. Mode~0
($\lambda=0.000$) tracks the invariant global ground state. Modes~1 and~2 exhibit exact
structural degeneracy ($\lambda=0.675$), manifesting as linear, mutually orthogonal spatial
gradients that function as coordinate charts over an isotropic 2D manifold. Higher modes
define localised geometric harmonics, with a UV cutoff set by the lattice spacing.}
\label{fig:eigenmodes}
\end{figure}

% ─────────────────────────────────────────────────────────────────────────────
\section{Spectral Foundations of the Substrate Energy Density}
\label{sec:spectral}
% ─────────────────────────────────────────────────────────────────────────────

\subsection{Canonical Partition Function and $\rho_0$}
The effective macroscopic energy density of the substrate is derived from
first-principles statistical mechanics. Treating the MI-Laplacian eigenvalues
$\{\lambda_n\}$ as the fundamental spectrum of structural fluctuations, we
define the canonical partition function over excited modes ($\lambda_n > 0$) as
\begin{equation}
Z = \sum_{n=1}^{N-1} e^{-\beta \lambda_n},
\label{eq:Z_def}
\end{equation}
where $\beta$ is an inverse entanglement temperature setting the
renormalisation-group coarse-graining scale.

\paragraph*{Entanglement temperature.}
The parameter $\beta$ in Eq.~\eqref{eq:Z_def} is the \emph{entanglement
temperature}: the Lagrange multiplier that fixes the average MI--Laplacian
excitation energy in the canonical ensemble of relational degrees of freedom.
In complete analogy with statistical mechanics, $\beta$ is conjugate to the
spectral ``energy'' $\lambda_n$ and determines the degree of coarse--graining
applied to the substrate. The ensemble average
\begin{equation}
\langle \lambda \rangle_\beta
= \frac{\sum_{n>0} \lambda_n e^{-\beta \lambda_n}}
       {\sum_{n>0} e^{-\beta \lambda_n}}
\label{eq:lambda_avg}
\end{equation}
is therefore fixed by the choice of $\beta$: large $\beta$ suppresses
high--eigenvalue modes and yields a smooth, long-wavelength geometric
description, while small $\beta$ retains short-wavelength fluctuations and
produces a more microscopic, rough geometry. The entanglement temperature thus
encodes the renormalisation-group coarse-graining scale of the emergent
geometry.

The internal energy density is then
\begin{equation}
\rho_{\rm eff}
= \frac{1}{V} \frac{\partial}{\partial \beta} \ln Z
= \frac{1}{V}
\frac{\sum_{n>0} \lambda_n e^{-\beta \lambda_n}}
     {\sum_{n>0} e^{-\beta \lambda_n}}.
\label{eq:rho_eff_def}
\end{equation}


\begin{equation}
\rho_0 = \int_{\lambda_{\min}}^{\lambda_{\max}} \lambda\,g(\lambda)\,d\lambda,
\label{eq:rho0}
\end{equation}
where $g(\lambda)$ is the MI-Laplacian spectral density, $\lambda_{\max}$ is set by the
correlation length at the transition, and $\lambda_{\min}$ is set by the largest connected
component. Crucially, $\rho_0$ is a spectral invariant of the relational graph: it is
determined entirely by the topology and connectivity of the substrate at the moment the
entanglement decay exponent collapses to $\alpha=6$.

Numerical evaluation on an $8\times8$ toroidal lattice at $\alpha=6$ yields
\begin{equation}
\rho_0 = 4.728\text{ (lattice units)},\quad
\lambda\in[0.834,\,8.148],\quad N_{\rm modes}=63.
\end{equation}
The Weyl-weighted spectral centroid (with $d_{\rm eff}=8$ measure) gives
$\rho_0^{\rm Weyl}=5.31$ lattice units, consistent with the spectral centroid to within
12\%.

\subsection{Physical Amplitude and the Role of $\varepsilon_0$}

Converting $\rho_0$ to physical units requires a microscopic energy scale
$\varepsilon_0$ (energy per eigenvalue unit) and an effective lattice volume $V$:
\begin{equation}
\rho_{\rm phys} = \frac{\rho_0\,N\,\varepsilon_0}{V}.
\end{equation}
The shape, tilt, UV behaviour, and loop finiteness of the resulting gravitational-wave
spectrum are fixed predictions of the substrate dynamics independent of $\varepsilon_0$.
The overall amplitude, however, is controlled by this microscopic energy scale. Systematic
exploration across energy scales from toy values to the Planck scale establishes that
physically plausible choices of $\varepsilon_0$ can reproduce PTA-range amplitudes, while
the spectral tilt $n_T\simeq1.875$ and UV cutoff remain unchanged. This situation is
analogous to inflationary cosmology, where the tilt of the scalar power spectrum is a
robust prediction while the amplitude depends on the inflaton potential height. We therefore
identify $\varepsilon_0$ as the sole free parameter of the amplitude sector and emphasise
that the parameter-independent predictions of the framework are the blue tilt and the
spectral structure of the gravitational-wave background.

\subsection{Two Distinct Dimension Concepts}
\label{sec:dims}

A source of potential confusion in the literature on emergent geometry is the conflation
of two distinct notions of dimension. We establish them rigorously here.

The \emph{Weyl spectral dimension} $d_{\rm eff}^{\rm Weyl}$ is extracted from the
MI-Laplacian eigenvalue counting function via
\begin{equation}
\lambda_n \propto n^{2/d_{\rm eff}^{\rm Weyl}},
\end{equation}
fitting the slope of $\log\lambda_n$ versus $\log n$ over the low-lying spectrum.

The \emph{fluid dimension} $d_{\rm fluid}$ is inferred from the equation-of-state scaling
\begin{equation}
\rho_{\rm eff}(a)\propto a^{-3(1+w_{\rm eff})},\quad
d_{\rm fluid} = 3(1+w_{\rm eff}) = \alpha,
\end{equation}
where the second equality uses $w_{\rm eff}=\alpha/3-1$ from the partition function
analysis \cite{JarvisInfoSat2026}.

These two quantities are \emph{not} equal in general. Numerical scans on the $6\times6$
toroidal lattice give the results summarised in Table~\ref{tab:dims}.

\begin{table}[h]
\centering
\caption{Weyl spectral dimension $d_{\rm eff}^{\rm Weyl}$ versus fluid dimension
$d_{\rm fluid}=\alpha$ at selected values of the MI decay exponent $\alpha$.}
\label{tab:dims}
\begin{ruledtabular}
\begin{tabular}{cccc}
$\alpha$ & $w_{\rm eff}$ & $d_{\rm eff}^{\rm Weyl}$ & $d_{\rm fluid}$\\
\hline
1.5 & $-0.500$ & 11.14 & 1.5\\
2.5 & $-0.167$ & 6.68  & 2.5\\
3.5 & $+0.167$ & 5.01  & 3.5\\
4.5 & $+0.500$ & 4.26  & 4.5\\
6.0 & $+1.000$ & 3.81  & 6.0\\
7.0 & $+1.333$ & 3.70  & 7.0\\
\end{tabular}
\end{ruledtabular}
\end{table}

At the hyper-stiff transition ($\alpha=6$, $w_{\rm eff}=1$) the Weyl spectral dimension
is $d_{\rm eff}^{\rm Weyl}=3.81$, not 8. The value $d_{\rm eff}\approx8$ reported in
earlier work \cite{JarvisEEG2026} corresponds to the critical power-law phase at
$\alpha=1.5$ ($w_{\rm eff}=-0.5$), which is a distinct regime. These two regimes must not
be conflated: the hyper-stiff amplification mechanism operates at $\alpha=6$, while the
hyper-dimensional geometry emerges at $\alpha=1.5$.

% ─────────────────────────────────────────────────────────────────────────────
\section{Holographic Saturation and Confinement Potential}
% ─────────────────────────────────────────────────────────────────────────────

\subsection{Information Bound}
A collapsing region of radius $r$ has an information storage capacity strictly bounded by
its boundary area according to the holographic principle \cite{tHooft1993, Susskind1995}:
\begin{equation}
I_{\max}(r) = \frac{A}{4} = \pi r^2,
\end{equation}
reaching complete information saturation at a limiting minimal radius $r=r_0$.

\subsection{Sign Reversal of $V_{\rm conf}$}
We introduce an effective macroscopic confinement potential $V_{\rm conf}(\rho)$:
\begin{equation}
V_{\rm conf}(\rho) =
\begin{cases}
0, & \rho < \rho_{\rm sat},\\
\leq 0, & \rho = \rho_{\rm sat},\\
< 0, & \rho > \rho_{\rm sat}.
\end{cases}
\end{equation}
This sign flip is the macroscopic thermodynamic signature of an underlying microphysical
phase transition, where the conventional attractive gravitational degrees of freedom
condense and manifest net repulsion at extreme densities.

% ─────────────────────────────────────────────────────────────────────────────
\section{GR-Consistent Junction Formulation}
% ─────────────────────────────────────────────────────────────────────────────

\subsection{Causal Boundary}
The collapse--expansion bounce is modelled entirely within classical GR across a spacelike
hypersurface $\Sigma$, which behaves as a strict localised causal boundary.

\subsection{Israel--Darmois Junction}
Matching the interior collapsing and exterior expanding spacetimes requires satisfying the
Israel--Darmois junction conditions. The induced metric $h_{ab}$ must be continuous:
\begin{equation}
[h_{ab}] = h_{ab}^+ - h_{ab}^- = 0.
\end{equation}
The jump in extrinsic curvature $K_{ab}$ is related to the surface stress-energy $S_{ab}$:
\begin{equation}
[K_{ab}] - h_{ab}[K] = -8\pi G\,S_{ab}.
\end{equation}
Because the transition occurs at global holographic saturation, $S_{ab}\to0$, and the
bounce is driven by a geometric momentum reflection across $\Sigma$ rather than exotic
matter fields.

\subsection{Modified Friedmann Equation}
Near the saturation hypersurface the effective dynamics are governed by
\begin{equation}
H^2 = \frac{\rho}{3}\left[1 - \chi(\rho)\right],
\end{equation}
with the smooth saturation envelope
\begin{equation}
\chi(\rho) = \frac{1}{2}\left[1+\tanh\!\left(\frac{\rho-\rho_{\rm sat}}{\Delta\rho}\right)\right].
\end{equation}
This guarantees a continuous, smooth variation in the cosmic deceleration parameters,
matching the continuous thermodynamic dilution of fields near the throat.

% ─────────────────────────────────────────────────────────────────────────────
\section{Entropy Inheritance and Renormalisation}
% ─────────────────────────────────────────────────────────────────────────────

In a cyclic universe, the absolute entropy produced during the expansion phase of cycle $n$,
combined with the residual entropy inherited from cycle $n-1$, dictates the total entropy
dumped into the subsequent bounce throat:
\begin{equation}
S_{\rm dump}^{(n)} = S_{\rm prod}^{(n)} + S_{\rm inh}^{(n)}.
\end{equation}
The dimensionless renormalised inhomogeneity ratio is
\begin{equation}
\tilde{S}^{(n+1)} = \frac{S_{\rm inh}^{(n+1)}}{S_{\max}^{(n+1)}}.
\end{equation}
While the absolute entropy $S_{\rm tot}$ increases monotonically in compliance with the
generalised second law, the total available phase space $S_{\max}$ expands at a faster
geometric rate across each bounce. The universe therefore re-enters the expansion phase in a
highly ordered, effectively low-entropy configuration without requiring non-unitary entropy
destruction.

% ─────────────────────────────────────────────────────────────────────────────
\section{Homogenisation}
% ─────────────────────────────────────────────────────────────────────────────

The generation of an isotropic background from an initially asymmetrical collapse proceeds
in two stages. During the late stages of the contracting phase, the extreme non-linearities
of the gravitational potential trigger widespread chaotic mixing, fragmenting large-scale
spatial anisotropies. As the density approaches $\rho_{\rm sat}$, holographic projection at
the boundary coarse-grains the remaining bulk fluctuations into $\mathcal{O}(A/\ell_P^2)$
organised boundary degrees of freedom, providing the uniform, isotropic initialisation
conditions required for standard FLRW expansion.

% ─────────────────────────────────────────────────────────────────────────────
\section{Cyclic Attractor}
% ─────────────────────────────────────────────────────────────────────────────

Each discrete cosmic bounce defines a discrete renormalisation-group-like flow
$\mathcal{R}$ on the space of effective background couplings:
\begin{equation}
\mathcal{R}:\{\lambda^{(n)}\}\mapsto\{\lambda^{(n+1)}\}.
\end{equation}

\begin{figure}[t]
\includegraphics[width=\linewidth]{entropy_attractor.png}
\caption{The Cosmic Renormalisation Attractor. The left panel shows the monotonic growth of
total entropy $S_{\rm tot}$ bounded safely within the growing capacity $S_{\max}$. The
right panel illustrates the renormalised inhomogeneity ratio $\tilde{S}(n)=S_{\rm
inh}/S_{\max}$ decaying exponentially to a stable fixed point at zero. The suppression
arises from a MI-Laplacian damping factor $e^{-0.22\,N/\rho_{\rm throat}}$ evaluated at
each throat, providing a first-principles mechanism for the inhomogeneity dilution.}
\label{fig:entropy_attractor}
\end{figure}

Numerical simulations of this flow establish the existence of a highly stable geometric
fixed point, as illustrated in Fig.~\ref{fig:entropy_attractor}. Even when initialised with
high structural asymmetry ($\tilde{S}^{(0)}=0.30$), the renormalised inhomogeneity ratio
$\tilde{S}(n)$ decays exponentially. By cycle~10 the system has converged to the stable
attractor $\tilde{S}^*=0$.

The damping factor in the simulation,
\begin{equation}
\mathcal{D}_n = \exp\!\left(-\frac{0.22\,N}{\rho_{\rm throat}^{(n)}}\right),
\end{equation}
is derived directly from the MI-Laplacian partition function: the ratio $N/\rho_{\rm
throat}$ measures how far the substrate energy density is from its saturation ceiling
$\rho_{\rm eff}^{\max}=N$, and the exponential suppression reflects the collapse of the
available phase space near holographic saturation. This provides a first-principles
derivation of the inhomogeneity suppression that was previously introduced
phenomenologically.

This convergence provides a robust, inflation-free solution to the Flatness and Horizon
problems, demonstrating that a homogeneous and flat cosmic geometry is the inevitable global
attractor of the cyclic network.

% ─────────────────────────────────────────────────────────────────────────────
\section{Scalar Perturbations Across the Holographic Saturation Bounce}
% ─────────────────────────────────────────────────────────────────────────────

\subsection{Perturbation Framework}
We evaluate the structural stability of the bounce by tracking linear cosmological
perturbations in the longitudinal gauge:
\begin{equation}
ds^2 = -(1+2\Phi)\,dt^2 + a(t)^2(1-2\Phi)\,d\vec{x}^2,
\end{equation}
where $\Phi$ denotes the gauge-invariant gravitational potential. The comoving curvature
perturbation $\zeta$ is
\begin{equation}
\zeta = -\Phi - \frac{H}{\dot{\rho}}\,\delta\rho.
\end{equation}
Transforming to conformal time, the Mukhanov--Sasaki variable $v=z\zeta$ satisfies
\begin{equation}
v_k'' + \left(c_s^2 k^2 - \frac{z''}{z}\right)v_k = 0.
\end{equation}

\subsection{Effective Sound Speed}
From the modified Friedmann framework,
\begin{equation}
\rho_{\rm eff} = \rho[1-\chi(\rho)],\qquad
\dot{H} = -\tfrac{1}{2}(\rho_{\rm eff}+p_{\rm eff}).
\end{equation}
The effective sound speed
\begin{equation}
c_s^2 = \frac{\partial p_{\rm eff}}{\partial\rho}
\end{equation}
remains bounded and strictly positive ($0<c_s^2\leq1$), ensuring the system stays clear
of gradient instabilities. At the hyper-stiff transition ($w_{\rm eff}=1$, $\alpha=6$) the
sound speed saturates its causal bound $c_s^2=w_{\rm eff}=1$, yielding $c_s=c$. This is
in sharp contrast to the earlier hand-tuned $w_{\rm conf}\simeq80$ parameterisation, which
implied $c_s\approx8.94c$ and violated causality. The derivation of $w_{\rm eff}=1$ from
first principles via the MI-Laplacian partition function therefore simultaneously fixes the
causal structure of the perturbation sector.

\subsection{Behaviour of $z''/z$}
As the universe passes through the bounce throat, $H\to0$. Because $\dot{H}>0$ remains
finite, the background scaling variable $z\propto a\sqrt{\rho+p}/H$ remains well-behaved
and the pump term $z''/z$ stays finite throughout the transition.

\subsection{Matching Conditions}
Applying the Israel--Darmois matching conditions across $\Sigma$ guarantees smooth
transmission of fluctuations:
\begin{equation}
[\Phi] = 0,\qquad[\zeta] = 0.
\end{equation}
Long-wavelength modes ($k\ll aH$) satisfy $\zeta'\approx0$, conserving primordial
amplitudes across the throat. Short-wavelength modes undergo stable adiabatic oscillations.
The bounce is therefore perturbatively stable with no ghost states or tachyonic
instabilities.

% ─────────────────────────────────────────────────────────────────────────────
\section{Phenomenology: Hyper-Stiff Phase and Tensor Spectrum}
\label{sec:pheno}
% ─────────────────────────────────────────────────────────────────────────────

\subsection{Confinement Energy and Hyper-Stiff Spike}

Residual higher-dimensional confinement energy left over from the geometric reflection
phase produces a lingering energy density component in the early expansion era:
\begin{equation}
\rho_{\rm conf}(a) = \rho_0\,\chi(a)[1-\chi(a)]^8,
\end{equation}
where $\chi(a)=\frac{1}{2}[1+\tanh((a_c-a)/\Delta a)]$. This generates a transient
hyper-stiff equation of state ($p>\rho$), with the peak value determined by the
MI-Laplacian partition function at $\alpha=6$:
\begin{equation}
w_{\rm eff} = \frac{\alpha}{3} - 1 = 1.0,\qquad
\rho_{\rm eff}(a)\propto a^{-6}.
\end{equation}
The resulting conformal potential
\begin{equation}
\frac{a''}{a} = a^2H^2(1-3w_{\rm eff}) = -2a^2H^2
\end{equation}
is strictly negative, amplifying sub-horizon tensor modes without violating causality
($c_s=c$) and without introducing ghost or gradient instabilities ($\omega_k^2=k^2+2a^2H^2>0$).
This resolves the earlier inconsistency in which a hand-tuned $w_{\rm conf}\simeq80$
produced a conformal potential $\approx-240a^2H^2$ and an acausal sound speed
$c_s\approx8.94c$.

\subsection{Tensor Spectral Index}

The tensor spectral index for a transient hyper-stiff phase with peak equation of state
$w_{\rm eff}$ and spike duration $\delta$ (measured in $\ln a$) is obtained by full
numerical integration of the Mukhanov--Sasaki equation
\begin{equation}
v_k'' + \omega_k^2(\tau)\,v_k = 0,\qquad
\omega_k^2 = k^2 - \frac{a''}{a},
\end{equation}
through the transient. The analytic formula
\begin{equation}
n_T^{\rm analytic} = \frac{6w_{\rm eff}-2}{1+3w_{\rm eff}}
\label{eq:nT_analytic}
\end{equation}
applies only in the limit of an eternal constant-$w$ background. For a finite-duration
Gaussian spike with $w_{\rm peak}=1$ this formula gives $n_T^{\rm analytic}=1.0$, which
underestimates the true tilt.

Full numerical integration of the Mukhanov--Sasaki equation across the transient, using an
implicit Radau solver with relative tolerance $10^{-8}$, yields
\begin{equation}
n_T = 1.875\qquad(\delta=0.5\text{ in }\ln a),
\end{equation}
with the result stable to $|\Delta n_T|\lesssim0.05$ over $\delta\in[0.4,\,0.6]$. The
discrepancy between the analytic and numerical values arises because, for a finite-duration
transient, modes that enter and exit the hyper-stiff window accumulate additional phase
relative to the eternal-background prediction. The numerically determined $n_T=1.875$ is
the physically correct prediction of the framework and supersedes the analytic estimate.

The tensor transfer function $\mathcal{T}(k)$ satisfies the adiabatic unitary bound
$\mathcal{T}(k)\leq1$ across the entire ultraviolet grid $k\in[10^2,10^4]$, confirming
loop finiteness:
\begin{equation}
\Delta^2_{\rm loop}\propto\int dk\,\mathcal{T}(k)\,k^3 < \infty.
\end{equation}

\subsection{Gravitational-Wave Amplitude}

The present-day spectral energy density is
\begin{equation}
\Omega_{\rm GW}(f) = \frac{1}{\rho_c}\frac{d\rho_{\rm GW}}{d\ln f}
\propto \left(\frac{f}{f_{\rm ref}}\right)^{n_T}
\end{equation}
in the blue-tilted band, with a physical turnover at the millihertz scale set by the finite
rupture duration. The overall amplitude is
\begin{equation}
\Omega_{\rm GW} = \frac{1}{12}\left(\frac{k}{a_0H_0}\right)^2
\Omega_{\rm rad}\,P_h,
\end{equation}
where the primordial tensor power $P_h=(2/\pi^2)(H_{\rm MI}/M_{\rm Pl})^2$ depends on
the Hubble rate during the hyper-stiff phase, which is in turn set by $\rho_{\rm phys}$.
For a substrate energy scale $\varepsilon_0$ near the Planck scale and a lattice volume
$V=N\ell_P^3$, the model yields $\Omega_{\rm GW}(f_{\rm PTA})$ in the PTA-observable range.
The spectral tilt $n_T=1.875$ and the UV behaviour of $\mathcal{T}(k)$ are
parameter-independent predictions; the amplitude is the sole quantity sensitive to
$\varepsilon_0$.

% ─────────────────────────────────────────────────────────────────────────────
\section{Discussion}
% ─────────────────────────────────────────────────────────────────────────────

The framework developed here differs from loop quantum cosmology \cite{AshtekarBojowald2006},
conformal cyclic cosmology \cite{PenroseBook}, and traditional fuzzball constructions
\cite{Susskind1995}. Rather than modifying the Einstein--Hilbert action with
higher-derivative terms or quantizing individual classical metric elements, it relies on a
GR-consistent junction driven by a holographic boundary condition.

The new spectral results of Section~\ref{sec:spectral} place the framework on firmer
quantitative footing. The identification of $\rho_0$ as a spectral invariant of the
MI-Laplacian (Eq.~\ref{eq:rho0}) removes the previously unconstrained normalisation
freedom in the confinement energy. The rigorous separation of $d_{\rm eff}^{\rm Weyl}$ and
$d_{\rm fluid}$ (Table~\ref{tab:dims}) eliminates an internal inconsistency present in
earlier drafts. The resolution of the $n_T$ discrepancy — from the analytic value of 1.0
to the numerically correct 1.875 — provides a sharp, falsifiable prediction for joint
PTA--LISA observations. The causality restoration ($c_s=c$ at $w_{\rm eff}=1$) removes
the most significant theoretical vulnerability of the mechanism.

Outstanding tasks include a full Boltzmann-code implementation to assess CMB and BAO
constraints, a derivation of $\varepsilon_0$ from holographic entropy bounds, and an
extension of the MI-Laplacian analysis to higher-dimensional substrates to test the
universality of the dimensional-deflation mechanism.

% ─────────────────────────────────────────────────────────────────────────────
\section{Conclusion}
% ─────────────────────────────────────────────────────────────────────────────

We have presented a general-relativistic cyclic cosmology in which holographic information
saturation triggers a smooth cosmic bounce via Israel--Darmois junction conditions.
Primordial scalar perturbations pass through the transition cleanly. New results establish:
(i) $\rho_0$ as a spectral invariant of the MI-Laplacian at $\alpha=6$; (ii) $n_T=1.875$
from full Mukhanov--Sasaki integration through the finite-duration transient, superseding
the analytic eternal-$w$ estimate of 1.0; (iii) a rigorous distinction between the Weyl
spectral dimension and the fluid equation-of-state dimension; (iv) full causality
($c_s=c$) at $w_{\rm eff}=1$; and (v) a first-principles derivation of the cyclic
attractor damping factor from the MI-Laplacian partition function. The GW amplitude
depends on the microscopic energy scale $\varepsilon_0$; the tilt and UV behaviour are
parameter-independent predictions testable by LISA, DECIGO, and multi-band PTA observations.

\begin{acknowledgements}
The author thanks the independent physics research community for structural feedback on
substrate dynamics scripts.
\end{acknowledgements}

\begin{thebibliography}{9}

\bibitem{Penrose1965}
R.~Penrose,
``Gravitational collapse and space-time singularities,''
Phys.\ Rev.\ Lett.\ \textbf{14}, 57 (1965).

\bibitem{HawkingPenrose1970}
S.~W.~Hawking and R.~Penrose,
``The singularities of gravitational collapse and cosmology,''
Proc.\ R.\ Soc.\ Lond.\ A \textbf{314}, 529 (1970).

\bibitem{AshtekarBojowald2006}
A.~Ashtekar and M.~Bojowald,
``Quantum geometry and the Schwarzschild singularity,''
Class.\ Quant.\ Grav.\ \textbf{23}, 391 (2006).

\bibitem{PenroseBook}
R.~Penrose,
\textit{Cycles of Time: An Extraordinary New View of the Universe}
(Bodley Head, London, 2010).

\bibitem{SteinhardtTurok2002}
P.~J.~Steinhardt and N.~Turok,
``Cosmic evolution in a cyclic universe,''
Science \textbf{296}, 1436 (2002).

\bibitem{tHooft1993}
G.~'t~Hooft,
``Dimensional reduction in quantum gravity,''
arXiv:gr-qc/9310026.

\bibitem{Susskind1995}
L.~Susskind,
``The world as a hologram,''
J.\ Math.\ Phys.\ \textbf{36}, 6377 (1995).

\bibitem{VanRaamsdonk2010}
M.~Van Raamsdonk,
``Building up spacetime with quantum entanglement,''
Gen.\ Rel.\ Grav.\ \textbf{42}, 2323 (2010).

\bibitem{JarvisCode}
P.~Jarvis,
``Emergent Geometry from Entanglement: A Numerical Demonstration using Heisenberg
Spin Chains'' (2026),
\url{https://github.com/ZaPpi3/emergent-geometry-heisenberg}.

\bibitem{JarvisEEG2026}
P.~Jarvis,
``Emergent Entanglement Geometry, Dynamical Dimensional Deflation, and Cosmological
Expansion Constraints'' (2026),
\url{https://github.com/ZaPpi3/hyperstiff-entanglement-geometry}.

\bibitem{JarvisGW2026}
P.~Jarvis,
``Blue-Tilted Primordial Gravitational Waves from a Transient Hyper-Stiff
Brane-Rupture Phase'' (2026),
\url{https://github.com/ZaPpi3/gravitational-wave-blue-tilt}.

\bibitem{JarvisInfoSat2026}
P.~Jarvis,
``The Geometry of Information Saturation: Mapping the Complete Cosmic Expansion
History from the Matrix Partition Function'' (2026),
\url{https://github.com/ZaPpi3/Information-saturated-bounce}.

\end{thebibliography}

\end{document}